\documentclass[11pt,a4paper]{article}
\usepackage[T1]{fontenc}
\usepackage{lmodern,microtype}
\usepackage[margin=25mm]{geometry}
\usepackage{amsmath,amssymb,amsthm,mathtools}
\usepackage{booktabs,enumitem,xcolor,fancyhdr}
\usepackage[hypertexnames=false,colorlinks=true,linkcolor=blue!45!black,urlcolor=blue!45!black]{hyperref}
\setlength{\emergencystretch}{3em}
\setlength{\parskip}{0.35em}
\setlist[itemize]{topsep=3pt,itemsep=3pt}
\pagestyle{fancy}\fancyhf{}
\fancyhead[L]{\small Inverse arctangent and Stieltjes rigidity}
\fancyhead[R]{\small Research manuscript}
\fancyfoot[C]{\thepage}
\setlength{\headheight}{14pt}
\newtheorem{theorem}{Theorem}
\newtheorem{lemma}[theorem]{Lemma}
\newtheorem{corollary}[theorem]{Corollary}
\title{An Inverse-Arctangent Function Is Not Stieltjes\\[4pt]
\large Moment rigidity, bounded phases, and an exact Hankel certificate}
\author{Keio Research Notes}
\date{31 October 2026}
\hypersetup{pdftitle={An Inverse-Arctangent Function Is Not Stieltjes},pdfauthor={Keio Research Notes}}
\begin{document}\maketitle
\begin{abstract}
We give a negative answer to the Stieltjes possibility raised in Question 1
of Keady, Khajohnsaksumeth, and Wiwatanapataphee (2018). The positive inverse
of $z\mapsto\arctan(1/z)/z$ is not a Stieltjes function. The main proof
uses positivity alone: three asymptotic coefficients would force a
representing measure to have zero variance, hence to be a point mass,
whereas an exact strict inequality rules out the resulting rational function.
We extend the obstruction to bounded monotone phases with a nonnegative
quadratic coefficient. Independently, we exhibit a strictly negative
$4\times4$ normalized derivative Hankel determinant, with its sign certified
by rational polynomial arithmetic. The complete-monotonicity questions
about this inverse and its square are not settled here.
\end{abstract}
\section{The published question and the conclusion}
G. Keady, N. Khajohnsaksumeth, and B. Wiwatanapataphee,
\emph{On functions and inverses, both positive, decreasing and convex:
And Stieltjes functions}, Cogent Mathematics \& Statistics 5:1 (2018),
article 1477543, DOI 10.1080/25742558.2018.1477543, define
\[
 \varphi_1(z)=\frac{\arctan(1/z)}{z},\qquad
 \varphi_2(z)=\varphi_1(\sqrt z),\qquad z>0.
\]
Question 1, on printed page 2 of 16, asks:
\begin{quote}
Is the inverse of $\varphi_2$ completely monotone?
Is the inverse of $\varphi_1$ completely monotone
(and, perhaps even Stieltjes)?
\end{quote}
On printed page 13, Section 5.2, the authors specifically state that
they have no obstruction to $\mu=\varphi_1^{-1}$ being Stieltjes.
They already know that $\mu_{(2)}=\varphi_2^{-1}=\mu^2$ is not
Stieltjes. The result below settles only the parenthetical Stieltjes
possibility for $\mu$. It does not prove or disprove complete
monotonicity of either $\mu$ or $\mu^2$.

\begin{theorem}
The positive inverse $\mu=\varphi_1^{-1}$ is not a Stieltjes function.
\end{theorem}

\section{Rigidity from three asymptotic coefficients}
If a Stieltjes function $F$ satisfies
\[
 F(x)=\frac{c}{x+d}+o(x^{-3})\qquad(x\to\infty),
 \quad c>0,\quad d\ge0,
\]
then $F(x)=c/(x+d)$ for every $x>0$.
To prove this, write its Stieltjes representation, absorbing any pole
at zero into the positive measure. Since $F(x)\to0$, its additive
constant is zero, and thus
\[
 F(x)=\int_{[0,\infty)}\frac{d\rho(t)}{x+t}.
\]
The first of the following three limits, followed by monotone
convergence, gives $\rho([0,\infty))=c$. Having established this,
the second limit gives the finite first moment $cd$. With that
moment established, the third limit gives the finite second moment
$cd^2$:
\begin{align*}
 xF(x)&=\int\frac{x}{x+t}\,d\rho(t)\longrightarrow c,\\
 x^2\left(\frac c x-F(x)\right)
 &=\int\frac{xt}{x+t}\,d\rho(t)\longrightarrow cd,\\
 x^3\left(F(x)-\frac c x+\frac{cd}{x^2}\right)
 &=\int\frac{xt^2}{x+t}\,d\rho(t)\longrightarrow cd^2.
\end{align*}
Each integrand is nonnegative and increases to the corresponding
moment integrand. No finiteness of a higher moment has been assumed.
It follows that
\[
 \int(t-d)^2\,d\rho(t)=cd^2-2d(cd)+d^2c=0.
\]
Positivity forces $\rho=c\delta_d$, proving the lemma.

\section{The inverse-arctangent obstruction}
The derivative
\[
 \varphi_1'(z)
 =-\frac{\arctan(1/z)}{z^2}-\frac1{z(1+z^2)}<0
\]
and its endpoint limits show that $\varphi_1$ is a smooth decreasing
bijection of $(0,\infty)$ onto itself. Put $L=\pi/2$. The inverse
relation is
\[
 y\mu(y)=L-\arctan\mu(y),\qquad y>0.
\]
It gives $0<\mu(y)<L/y$. Since
\[
 0<u-\arctan u=\int_0^u\frac{t^2}{1+t^2}\,dt<\frac{u^3}{3}
 \qquad(u>0),
\]
we obtain the exact, strict, all-$y$ estimate
\begin{equation}
 0<\mu(y)-\frac{L}{y+1}
   =\frac{\mu(y)-\arctan\mu(y)}{y+1}
   <\frac{L^3}{3y^3(y+1)}. \tag{R}
\end{equation}
In particular $\mu(y)=L/(y+1)+O(y^{-4})$. If $\mu$ were
Stieltjes, the moment-rigidity lemma would force equality with
$L/(y+1)$ everywhere, contradicting the strict lower inequality
in (R). This proves the theorem without numerical computation or
any derivative beyond the one needed to define the inverse.

\subsection{Positive scale parameters}
For every $\alpha,b>0$, the positive inverse of
\[
 \varphi_{\alpha,b}(z)=\alpha\frac{\arctan(b/z)}{z}
\]
is not Stieltjes. Indeed it is $g_{\alpha,b}(y)=b\mu(by/\alpha)$,
and the Stieltjes class is invariant under positive multiplication
and positive scaling of the argument in either direction.
In the Robin normalization $\varphi_\beta(z)=\arctan(1/(\beta z))/z$
with $\beta>0$, this inverse is $\beta^{-1}\mu(y/\beta)$, so the
negative conclusion applies to every positive Robin parameter.
The limiting elementary inverse at $\beta=0$ is $L/y$, which is
Stieltjes.

\section{A necessary condition for bounded phases}
Let $L>0$ and let $q:[0,\infty)\to[0,L)$ be nonzero, with
$q(0)=0$, $q\in C^1(0,\infty)$, a $C^2$ extension at zero, and
$q'(z)\ge0$ for $z>0$. Define
\[
 \Phi(z)=\frac{L-q(z)}{z},\qquad z>0.
\]
Then $\Phi$ is a strictly decreasing bijection of $(0,\infty)$
onto itself. If its inverse $u=\Phi^{-1}$ is a Stieltjes function,
it is necessary that
\[
 \boxed{q''(0)<0.}
\]
In particular, $q''(0)\ge0$ rules out the Stieltjes property.

\emph{Proof.} The strict derivative inequality
\[
 \Phi'(z)=\frac{-zq'(z)-L+q(z)}{z^2}<0
\]
and the endpoint limits establish the bijection. Write
$q(z)=az+bz^2+o(z^2)$, where $a=q'(0)\ge0$ and
$b=q''(0)/2$. The inverse identity $yu(y)=L-q(u(y))$, initially
giving $0<u(y)\le L/y$, yields successively
\[
 u(y)=\frac L y-\frac{aL}{y^2}
              +\frac{a^2L-bL^2}{y^3}+o(y^{-3}).
\]
If $u$ is Stieltjes, exactly the three monotone-convergence steps
in the rigidity lemma give a finite positive representing measure
with moments
\[
 m_0=L,\qquad m_1=aL,\qquad m_2=a^2L-bL^2.
\]
Its squared deviation from $a$ is
\[
 \int(t-a)^2\,d\rho(t)=-bL^2\ge0,
\]
so $b\le0$. If $b=0$, the measure must be $L\delta_a$, hence
$u(y)=L/(y+a)$. Substitution into the inverse relation, using that
$u$ ranges over all of $(0,\infty)$, forces $q(z)=az$ for every
$z>0$. Boundedness of $q$ forces $a=0$, and then $q\equiv0$,
contrary to the hypotheses. Thus $b<0$.

The result applies, under the displayed monotonicity and smoothness
hypotheses, to every bounded nonzero odd phase: its second derivative
at zero vanishes. In particular it applies to $q(z)=\arctan z$
with $L=\pi/2$, and to $q(z)=\tanh z$ or $q(z)=\operatorname{erf}z$
with $L=1$. The result for these additional phases is a derived
extension; no publication-priority claim is made.

The strict necessary sign is nonvacuous. For example,
$q(z)=Lz/(1+z)$ has $q''(0)=-2L<0$, and its inverse is Stieltjes:
\[
 u(y)=\frac{\sqrt{1+4L/y}-1}{2}
 =\frac1{2\pi}\int_0^{4L}\frac1{y+t}
                         \sqrt{\frac{4L-t}{t}}\,dt.
\]
The equality follows by $t=4L\sin^2\theta$. This example does not
assert that the sign condition is sufficient for arbitrary phases.

\section{An independent finite Hankel certificate}
The following exact local certificate independently proves the same
theorem and records how a finite moment-Gram test detects the
obstruction. Set
\[
 a=\frac\pi4,\qquad \tau=\frac1{1+2a}=\frac2{2+\pi},\qquad
 s_n=\frac{(-1)^n}{n!}\mu^{(n)}(a).
\]
In particular $\mu(a)=1$. Differentiating
$y\mu(y)=\arctan(1/\mu(y))$ gives the exact differential identity
\begin{equation}
 [1+y(1+\mu^2)]\mu'+\mu+\mu^3=0. \tag{A}
\end{equation}
Substitution of the Taylor series at $a$ into (A) yields
\begin{align*}
 s_0={}&1,\\
 s_1={}&2\tau,\\
 s_2={}&4\tau^2+2\tau^3,\\
 s_3={}&8\tau^3+\frac{32}{3}\tau^4+4\tau^5,\\
 s_4={}&16\tau^4+\frac{112}{3}\tau^5+\frac{100}{3}\tau^6+10\tau^7,\\
 s_5={}&32\tau^5+\frac{1624}{15}\tau^6+\frac{488}{3}\tau^7
                +112\tau^8+28\tau^9,\\
 s_6={}&64\tau^6+\frac{1416}{5}\tau^7+\frac{9128}{15}\tau^8
                +\frac{6272}{9}\tau^9+392\tau^{10}+84\tau^{11}.
\end{align*}
These expressions can be checked independently, without (A), by
inverting the Taylor series of $(\pi/2-\arctan(1+h))/(1+h)$.
Its first six nonconstant coefficients, multiplied by $2\tau$, are
\[
 -1,\quad 1+\frac\tau2,\quad -1-\frac{2\tau}3,\quad
 1+\frac{2\tau}3,\quad -1-\frac{37\tau}{60},\quad
 1+\frac{23\tau}{40}.
\]
The companion exact-arithmetic script verifies this composition and
the following determinant identity as polynomial identities over
$\mathbb Q[\tau]$:
\begin{equation}
 \det(s_{i+j})_{i,j=0}^3
 =\frac8{2025}\tau^{15}P(\tau),\qquad
 P(t)=-96+1280t-4540t^2-2820t^3+8100t^4+8100t^5+2025t^6.
 \tag{B}
\end{equation}
Only $3<\pi<4$ is needed for its sign: this gives
$1/3<\tau<2/5$. Put $u=15\tau-5\in(0,1)$. Direct expansion gives
\[
 P\left(\frac{u+5}{15}\right)
 =-\frac{1279}{9}-\frac{562}{9}u+\frac{283}{45}u^2
   +\frac{1232}{225}u^3+\frac{37}{75}u^4
   +\frac2{125}u^5+\frac1{5625}u^6
 \le-\frac{730334}{5625}<0.
\]
Alternatively, in the Bernstein basis
$\binom6j u^j(1-u)^{6-j}$, the seven coefficients have common
denominator $16875$ and numerator vector
\[
 (-2398125,-2573750,-2742300,-2899155,
   -3039140,-3156480,-3244752),
\]
another exact certificate of strict negativity.

For any Stieltjes function $f$ and any $a>0$, its normalized signed
derivatives $m_n=(-1)^nf^{(n)}(a)/n!$ are moments of a finite
positive measure. Indeed, from
\[
 f(y)=c_0+\frac{c_1}{y}+\int_{[0,\infty)}\frac{d\sigma(t)}{y+t},
 \qquad c_0,c_1\ge0,\quad \sigma\ge0,
\]
the constant is an atom of mass $c_0$ at $0$, the pole is an atom
of mass $c_1/a$ at $1/a$, and the remaining measure is the pushforward
of $(a+t)^{-1}d\sigma(t)$ under $t\mapsto(a+t)^{-1}$.
All moments exist since the support is contained in $[0,1/a]$.
Thus every Hankel matrix $(m_{i+j})$ is positive semidefinite:
for a real vector $v$, its quadratic form is
$\int(\sum_i v_i r^i)^2d\nu(r)\ge0$.
The strictly negative determinant in (B) contradicts this necessary
condition for $\mu$ to be Stieltjes, providing a second proof.

\section{Source audit and exact limits of the claim}
The publisher's indexed PDF directly supplies Question 1 on page 2
and Section 5.2 on page 13. The HTML publisher record gives publication
on 14 June 2018 (the PDF also records the earlier accepted-version
date 29 May 2018).
Searches for the specific inverse-arctangent Stieltjes question,
including Keady plus Hankel/determinant terms and searches restricted
to the Pudim repository, found no earlier proof of the theorem above.
This is not an exhaustive publication-priority certification.

The public Pudim application ledger entry APP-0035 concerns the
different Question 3, about arbitrary completely monotone self-bijections.
It gives a counterexample of the form $x^{-1}+100e^{-x}$; that result
does not address $\mu=\varphi_1^{-1}$. The full rendered application
ledger was searched: it contains no occurrence of \texttt{arctan},
and its Keady entry concerns only Question 3. It does not claim the
specific Stieltjes result proved here. Its full
large TeX theory file could not be retrieved through the available
public-fetch tool, so no claim of an exhaustive full-repository audit
is made.

The 2025 Jovanovi\'c--Treml paper on the \emph{reciprocal}
$1/\arctan x$ concerns a different function; it is not a result
about the compositional inverse $\varphi_1^{-1}$.

\subsection*{References and source links}
\begin{itemize}
\item Publisher paper:
\href{https://www.tandfonline.com/doi/full/10.1080/25742558.2018.1477543}{Keady--Khajohnsaksumeth--Wiwatanapataphee (2018), publisher record}
\item Publisher PDF:
\href{https://www.tandfonline.com/doi/pdf/10.1080/25742558.2018.1477543}{Article 1477543, publisher PDF}
\item Related Robin-eigenvalue paper, definitions and context:
\url{https://files.ele-math.com/articles/mia-21-62.pdf}
\item Public competing-claim ledger:
\href{https://github.com/pudim-project/pudim-ai-demo-zetalaw/blob/main/APPLICATIONS.md}{Pudim application ledger (APP-0035, GitHub)}
\end{itemize}

\end{document}
